\documentclass[12pt]{beamer}
\usetheme{moloch}

\usepackage{amsmath,amscd}
\usepackage{fontspec}
\usepackage{unicode-math}
\usepackage[dvipsnames]{xcolor}

\setbeamercolor{background canvas}{bg=black}
\setbeamercolor{normal text}{fg=white}

\setmonofont{DejaVu Sans Mono}

% Set the title bar to black background with white text
\setbeamercolor{frametitle}{bg=black, fg=white}

% Ensure other structural elements (bullets, etc.) are visible
\setbeamercolor{structure}{fg=white}

\setmainfont{Libertinus Serif}
\setsansfont{Libertinus Sans}
\setmonofont{DejaVu Sans Mono}
\setmathfont{STIXTwoMath-Regular.otf}
%\setmathfont{STIX Two Math}


% Define colors for syntax highlighting
\definecolor{codegreen}{rgb}{0,0.6,0}
\definecolor{codegray}{rgb}{0.5,0.5,0.5}
\definecolor{codecoffee}{rgb}{0.75,0.6,0.3}
\definecolor{question}{rgb}{1.0,0.75,0.1}
\definecolor{lightblue}{rgb}{0.55,0.75,0.9}
\definecolor{yellow}{rgb}{1.0, 1.0, 0.0}
\definecolor{codecomment}{rgb}{0.97, 0.74, 0.4}
\definecolor{codepy}{rgb}{0.9, 0.85, 0.85} %{#F8BC65}
\definecolor{funcColor}{rgb}{0.47, 0.6, 0.8}
\definecolor{midgray}{rgb}{0.4, 0.4, 0.3}
\definecolor{indexwhite}{rgb}{0.9, 0.9, 0.9}

\newcommand{\FrameTitle}[2]{\frametitle{#1 \hfill \small\textnormal{\textcolor{midgray}{#2}}}}
\newcommand{\snl}{\vspace*{4pt}\\}
\newcommand{\nl}{\vspace*{1em}\\}
\newcommand{\vsp}{\vspace*{1em}}
\newcommand{\SixPtVspace}{\vspace*{6pt}}

\newcommand{\Quiz}[1]{{\large\textbf{\color{question}{#1}}}}
\newcommand{\QuizQuestion}[1]{{\Large\textbf{\color{question}{#1}}}}

\newcommand{\Matrix}[1]{\textbf{#1}}
\newcommand{\Vector}[1]{\textbf{#1}}
\newcommand{\MMat}[1]{\mathbf{#1}}
\newcommand{\MVec}[1]{\mathbf{#1}}
\newcommand{\Norm}[1]{\lVert {{#1}} \rVert}
\newcommand{\NormBeta}[1]{\lVert \beta \mathbf{#1} \rVert}
\newcommand{\NormConst}[1]{\lVert #1 \rVert}
\newcommand{\Trans}{\text{\raisebox{0.5ex}{\scalebox{0.75}{T}}}}
\newcommand{\abs}[1]{\lvert {#1} \rvert}
\newcommand{\Avg}[1]{\mathbf{avg}{(\MVec{#1})}}
\newcommand{\AvgVec}[1]{\mathbf{avg}{(\MVec{#1})}}
\newcommand{\AvgDemeanedVec}[1]{\mathbf{avg}{(\tilde\MVec{#1})}}
\newcommand{\DemeanedVec}[1]{\tilde\MVec{#1}}
\newcommand{\Std}[1]{\mathbf{std}(\MVec{#1})}

\begin{document}
\title{\center{Applied Linear Algebra\\\vsp Standard Deviation\\}}

\author{\center{Devi Prasad M \\ Manipal School of Information Sciences \\ Manipal \\}}
\date{\center{September 2026}}
\maketitle

\begin{frame}\FrameTitle{De-meaned Vector}{Unit 1}
    For any vector $\MVec{x} \in \mathbb{R}^n$, the \emph{mean} of its entries is
    \[
        \mathbf{avg}(\MVec{x}) = \frac{1}{n}\sum_{i=1}^{n} x_i
    \]
    and the associated \emph{de-meaned} vector is
    \[
        \DemeanedVec{x} = \MVec{x} - \AvgVec{x}\MVec{1}_n
    \]

    Notice $\AvgDemeanedVec{x} = 0$, that is, the mean
    value of the entries of $\tilde{\MVec{x}}$ is zero.
\end{frame}

\begin{frame}\FrameTitle{De-meaned Vector}{Unit 1}
    Let $\MVec{x} = (10, -2, 8, 23, 6, 75)$.

    \vsp
    Then $\Avg{x} = 20$ and the de-meaned vector is
    \begin{align*}
        \tilde{\MVec{x}} &= \MVec{(10, -2, 8, 23, 6, 75)} - \MVec{(20, 20, 20, 20, 20, 20)}\\[4pt]
                            &=(10-20, -2-20, 8-20, 23-20, 6-20, 75-20)\\[4pt]
                            &= (-10, -22, -12, 3, -14, 55)
    \end{align*}

    \vsp
    The mean of the entries of $\tilde{\MVec{x}}$ is zero.
\end{frame}

\begin{frame}\FrameTitle{De-meaned Vector}{Unit 1}
    \QuizQuestion{What is $\tilde{\MVec{x}}$ if all entries of $\MVec{x}$ are equal?}
\end{frame}

\begin{frame}\FrameTitle{Norm of a De-meaned Vector}{Unit 1}
    Recall $\MVec{x} = (10, -2, 8, 23, 6, 75)$ and $\Avg{x} = 20$. \vspace*{4pt}\\
    $\Norm{\MVec{x}} = \sqrt{100+4+64+529+36+5625} = \sqrt{6358} \approx 79.7371$.\vspace*{8pt}\\

    \noindent\rule{\textwidth}{0.4pt}
    \begin{align*}
    \tilde{\MVec{x}} &= \MVec{(10, -2, 8, 23, 6, 75)} - \MVec{(20, 20, 20, 20, 20, 20)}\\[2pt]
                            &= (-10, -22, -12, 3, -14, 55)
    \end{align*}

    Norm of the de-meaned vector $\tilde{\MVec{x}}$ is
    \begin{align*}
        \Norm{\tilde{\MVec{x}}} &= \sqrt{(-10)^2 + (-22)^2 + (-12)^2 + 3^2 + (-14)^2 + 55^2} \\[2pt]
                        &= \sqrt{10^2 + 22^2 + 12^2 + 3^2 + 14^2 + 55^2}\\[2pt]
                        &= \sqrt{100 + 484 + 144 + 9 + 196 + 3025} \\[2pt]
                        &= \sqrt{3958} \\[2pt]
                        &\approx 62.9126
    \end{align*}
\end{frame}

\begin{frame}\FrameTitle{Standard Deviation}{Unit 1}
    The \emph{standard deviation} of an $n$-vector $\MVec{x}$
    is defined as the RMS value of the de-meaned vector $\MVec{x} - \Avg{x}\MVec{1}_n$:

    \[
        \begin{array}{rll}
            \Std{\MVec{x}} &=& \sqrt{\dfrac{1}{n}\sum_{i=1}^{n} (x_i - \Avg{x})^2} \\[1.5em]
            &=& \sqrt{\dfrac{(x_1-\Avg{x})^2 + \cdots + (x_n-\Avg{x})^2}{n}} \\[1.5em]
            &=& \dfrac{\Norm{\tilde{\MVec{x}}}}{\sqrt{n}} \\[1.5em]
            &=& \dfrac{\lVert\,\MVec{x} - ((\MVec{1}_n^{\Trans}\,\MVec{x})/{n}) \, \MVec{1}_n \,\rVert}{\sqrt{n}}
        \end{array}
    \]
\end{frame}

\begin{frame}\FrameTitle{Standard Deviation}{Unit 1}
Consider our example
\begin{align*}
    \MVec{x} &= (10, -2, 8, 23, 6, 75)\\[4pt]
    n &= 6\\[4pt]
    \Avg{x} &= 20\\[4pt]
    \tilde{\MVec{x}} &= (-10, -22, -12, 3, -14, 55)\\[4pt]
    \Norm{\tilde{\MVec{x}}} &\approx 62.9126
\end{align*}

Now we can compute the standard deviation of $\MVec{x}$:
\begin{align*}
    \Std{\MVec{x}} &= \Norm{\tilde{\MVec{x}}}/\sqrt{n} \\[4pt]
                   &= 62.9126/\sqrt{6} \\[4pt]
                   &\approx 25.686
\end{align*}
\end{frame}

\begin{frame}\FrameTitle{Standard Deviation}{Unit 1}
Given two vectors $\MVec{x}$ and $\MVec{y}$,
where the value of each entry of $\MVec{y}$ is
$\Std{\MVec{x}}$,

\begin{align*}
    \MVec{x} &= (-10, -22, -12, 3, -14, 55), \ and \\[4pt]
    \MVec{y} &= (25.686, 25.686, 25.686, 25.686, 25.686, 25.686),
\end{align*}

\QuizQuestion{what is $\Norm{\MVec{x}}$ and $\Norm{\MVec{y}}$?}
\end{frame}

\begin{frame}\FrameTitle{Standard Deviation - A ``Good'' Sample}{Unit 1}
    \[
        \text{Let} \ \MVec{x} = (10, 2, 9, 18, 6, 15, 12, 7, 11)
    \]
    \begin{align*}
        \Avg{x} &= 10 \\[3pt]
        \Norm{\MVec{x}} &= \sqrt{100 + 4 + 81 + 324 + 36 + 225 + 144 + 49 + 121} \\[3pt]
                        &= \sqrt{1084} \approx 32.9242 \\[3pt]
        \tilde{\MVec{x}} &= (0, -8, -1, 8, -4, 5, 2, -3, 1) \\[3pt]
        \Norm{\tilde{\MVec{x}}} &= \sqrt{0+64+1+64+16+25+4+9+1}\\[3pt]
                                   &= \sqrt{184} \approx 13.5647\\[3pt]
        \Std{\MVec{x}} &= \dfrac{13.5647}{\sqrt{9}} = 4.522
    \end{align*}
\end{frame}

\begin{frame}\FrameTitle{Standard Deviation}{Unit 1}
    \begin{itemize}
    \item The standard deviation of a vector is \textbf{zero} only
            when all its entries are equal.
    \item The standard deviation of a vector is \emph{small} when the
            entries of the vector are nearly the same.
    \end{itemize}

    \vspace*{2em}
    The standard deviation of a vector $\MVec{x}$ is useful when
    the data set does not contain outliers, as it provides a more
    accurate measure of the spread of the data.
\end{frame}

\begin{frame}\FrameTitle{Properties of Standard Deviation}{Unit 1}
    Assume a vector $\MVec{x}$ and any number $a$.
    \begin{itemize}
        \item Adding a constant to every entry of a vector does not change its
              standard deviation. \vspace*{4pt}\\
            $\Std{\MVec{x} + {\mathit{a}}\,\MVec{1_{\mathit{n}}}} = \Std{\MVec{x}}$

        \item[]

        \item  Multiplying a vector by a scalar multiplies the standard
               deviation by the absolute value of the scalar.\vspace*{4pt}\\
            $\Std{\mathit{a}\,\MVec{x}} = \abs{a}\,\Std{\MVec{x}}$
    \end{itemize}
\end{frame}

\begin{frame}\FrameTitle{Standardization}{Unit 1}
The \emph{standardization} of a vector $\MVec{x}$ is the vector
\[
    \MVec{z} = \dfrac{\MVec{x} - \Avg{\MVec{x}}\MVec{1}_n}{\Std{\MVec{x}}}
             = \dfrac{\DemeanedVec{x}}{\Std{\MVec{x}}}
\]

\vsp
$\MVec{z}$ is the \emph{standardized} version of $\MVec{x}$, and has the following properties:
\begin{itemize}
    \item $\Avg{\MVec{z}} = 0$
    \item $\Std{\MVec{z}} = 1$
\end{itemize}

\vsp
The entries of $\MVec{z}$ are called the $z$-scores associated with the
original entries of $\MVec{x}$.
\end{frame}


\begin{frame}\FrameTitle{Correlation Coefficient}{Unit 1}
    Let $\MVec{a}$ and $\MVec{b}$ be two $n$-vectors, with the
    associated de-meaned vectors $\DemeanedVec{a}$ and $\DemeanedVec{b}$.
    \[
        \DemeanedVec{a} = \MVec{a} - \Avg{a}\mathbf{1},
        \qquad\qquad
        \DemeanedVec{b} = \MVec{b} - \Avg{b}\mathbf{1}
    \]

    The \emph{correlation coeficient} of $\DemeanedVec{a}$ and
    $\DemeanedVec{b}$ is defined as
    \[
        \rho = \dfrac{\DemeanedVec{a}^{\Trans} \DemeanedVec{b}}
        {\Norm{\DemeanedVec{a}} \Norm{\DemeanedVec{b}}}
    \]
    assuming that the de-meaned vectors are not $\mathbf{0}$.
\end{frame}

\begin{frame}\FrameTitle{Correlation Coefficient}{Unit 1}
    It is clear from the definition of $\rho$ that it is
    cos $\theta$, the angle between $\DemeanedVec{a}$ and
    $\DemeanedVec{b}$.
    \[
        \rho = \angle (\DemeanedVec{a}, \DemeanedVec{b})
    \]
\end{frame}


\begin{frame}\FrameTitle{Correlation Coefficient}{Unit 1}
    \begin{itemize}
        \item The correlation coefficient, $\rho$, tells us how the
                entries in the two vectors vary together.
        \item $-1 \le \rho \le 1$
        \item If $\rho = 0$, the vectors are $uncorrelated$.
        \item A vector with all entries equal is said to be
                \emph{uncorrelated} with any vector.
    \end{itemize}
\end{frame}


\begin{frame}\FrameTitle{Correlation Coefficient}{Unit 1}
    Define the standardized versions of $\MVec{a}$ and $\MVec{b}$.
    \[
        u = \DemeanedVec{a} / \Std{a}, \ and \ v = \DemeanedVec{b}/\Std{b}
    \]

    Then,
    \[
        \rho = {u^{\Trans} v}/{n}
    \]
\end{frame}

\begin{frame}\FrameTitle{Correlation Coefficient}{Unit 1}
\end{frame}

\begin{frame}\FrameTitle{Correlation Coefficient}{Unit 1}
\end{frame}


\end{document}